% ---------------------------------------------------------------------------
% Evaluation: data-parallel training of the residual corrector.
% Numbers: SZ3 block   = bench_out/table (NYX x3)
%                        + table_qmc_ep20 (QMCPack, 20 epochs to cover the 60 s budget)
%                        + table_rand (Magnetic, random sampling)
%                        + table_rand_bo (Miranda, random sampling + Phase-1 BO:
%                          axis 1, lr 3.12e-3)
%          SPERR block = bench_out/table_sperr (NYX x3)
%                        + table_sperr_qmc_ep20 (QMCPack: qmc_n1 at 20 epochs,
%                          qmc_n4_ep40 at 40 -- the 20-epoch schedule saturated
%                          0.13 dB under the target)
%                        + table_sperr_rand (Miranda lr 4e-3, Magnetic; random
%                          sampling)
% Regenerate with: python bench_compress_table.py --latex bench_out/table/nyx_*.json \
%   bench_out/table_qmc_ep20/*.json bench_out/table_rand/mag_*.json \
%   bench_out/table_rand_bo/*.json bench_out/table_sperr/nyx_*.json \
%   bench_out/table_sperr_qmc_ep20/qmc_n1.json \
%   bench_out/table_sperr_qmc_ep20/qmc_n4_ep40.json bench_out/table_sperr_rand/*.json
%   (default --iso budget; --iso psnr = min-final variant, --iso work = 10-epoch.
%    Superseded runs kept for the record: table_240k, table/{miranda,mag},
%    table/qmc (10 ep), table_sperr/{miranda,mag,qmc},
%    table_sperr_qmc_ep20/qmc_n4 (20-ep schedule, saturates under target).)
% Figure: psnr_vs_walltime_2panel_cr500.pdf (SZ3-NYX | SPERR-Miranda), from
%   plot_time_to_psnr.ipynb (run all cells). The retired single-panel figures
%   (psnr_vs_walltime_cr500.pdf, psnr_vs_walltime_sperr_cr500.pdf, both NYX)
%   are still regenerated by plot_time_to_psnr.py.
% See NOTES FOR REVISION at the end before submitting.
% ---------------------------------------------------------------------------

\subsection{Parallel Training Throughput}
\label{sec:eval-parallel}

The residual corrector is trained once per field at compression time, so its
cost is charged to the compressor. We measure how much of that cost four GPUs
remove.

\paragraph{Setup.} We compress six fields from four SDRBENCH datasets at
$\mathrm{CR}\!\approx\!500$ on four NVIDIA A100 80\,GB PCIe GPUs, using
\textsc{DistributedDataParallel} over NCCL, with the corrector trained on top
of each of the two base codecs: SZ3 and SPERR. Both configurations run 10
epochs and take the same number of slices per optimizer update, so they differ
only in GPU count. Table~\ref{tab:parallel-config} gives the per-dataset
settings. Compression time is codec encoding plus training; decompression time
is codec decoding plus one full-volume forward pass. Per-epoch evaluation is
measurement, not compression, so it is excluded from the training clock.

\begin{table}[t]
  \centering
  \caption{Per-dataset configuration. Batch is slices per optimizer update,
  written as single-GPU / per-rank: four GPUs split the same batch, so both see
  the same data per update. QMCPack takes 256 because its slices are
  $55\times$ smaller than NYX's. Slices are drawn randomly for Miranda,
  QMCPack, and Magnetic Reconnection and sequentially for NYX; Miranda slices
  along axis~1 with lr $3.1\times10^{-3}$ (both chosen by the Phase-1 search),
  the others along axis~0. $\mathrm{CR}$ is effective, so model weights count
  toward the payload.}
  \label{tab:parallel-config}
  \begin{tabular}{lrrrrr}
    \toprule
    Dataset & Shape & Size & Batch & Params & CR \\
            &       & (GB) & (1\,/\,4\,GPU) & & \\
    \midrule
    NYX -- baryon density      & $512^3$               & 0.50 & 4 / 1    &  29\,803 & 445 \\
    NYX -- temperature         & $512^3$               & 0.50 & 4 / 1    &  29\,803 & 452 \\
    NYX -- dark matter density & $512^3$               & 0.50 & 4 / 1    &  29\,803 & 450 \\
    Miranda                    & $1024^3$              & 4.00 & 4 / 1    & 237\,538 & 450 \\
    QMCPack                    & $33120\!\times\!69^2$ & 0.59 & 256 / 64 &  28\,858 & 458 \\
    Magnetic Reconnection      & $512^3$               & 0.50 & 4 / 1    &  28\,858 & 451 \\
    \bottomrule
  \end{tabular}
\end{table}

\begin{table}[t]
  \centering
  \caption{End-to-end compression and decompression cost at
  $\mathrm{CR}\!\approx\!500$ for both base codecs on one versus four GPUs,
  at equal quality. Each row's target PSNR is what the single GPU attains
  within the training budget the method is configured with (10\,s for the
  $512^3$ fields, 60\,s for QMCPack, 80\,s for Miranda); compression is codec
  encoding plus residual-model training to that target, decompression is
  codec decoding plus one full-volume forward pass. Both configurations
  perform the same number of slices per optimizer update, so within a row
  they differ only in GPU count; the speed-up is the ratio of training times.
  SZ3 is timed in-process, SPERR through its command-line tool.}
  \label{tab:parallel-time}
  \begin{tabular}{lrrrrrr}
    \toprule
    & & \multicolumn{2}{c}{Single GPU} & \multicolumn{2}{c}{4 GPUs} & \\
    \cmidrule(lr){3-4}\cmidrule(lr){5-6}
    Dataset & PSNR & Comp. & Dec. & Comp. & Dec. & Speed-up \\
            & (dB) & (s)   & (s)  & (s)   & (s)  &          \\
    \midrule
    \multicolumn{7}{l}{\textit{SZ3 + Ours}} \\
    NYX -- baryon density      & 121.67 &  11.59 &  1.96 &  4.52 &  1.92 & 3.48$\times$ \\
    NYX -- temperature         &  83.40 &  11.58 &  1.93 &  5.20 &  1.90 & 2.82$\times$ \\
    NYX -- dark matter density &  84.19 &  11.55 &  1.83 &  5.11 &  1.79 & 2.82$\times$ \\
    Miranda                    &  44.69 &  94.72 & 25.49 & 35.16 & 24.70 & 3.80$\times$ \\
    QMCPack                    &  65.10 &  61.82 &  2.80 & 23.09 &  2.84 & 2.82$\times$ \\
    Magnetic Reconnection      &  50.48 &  11.75 &  1.66 &  5.29 &  1.65 & 2.87$\times$ \\
    \midrule
    \multicolumn{7}{l}{\textit{SPERR + Ours}} \\
    NYX -- baryon density      & 113.25 &  10.94 &  2.09 &  4.00 &  2.11 & 3.30$\times$ \\
    NYX -- temperature         &  75.90 &  10.92 &  2.13 &  3.39 &  2.08 & 4.08$\times$ \\
    NYX -- dark matter density &  80.95 &  10.95 &  2.14 &  3.92 &  2.08 & 3.41$\times$ \\
    Miranda                    &  49.24 &  87.01 & 28.87 & 38.85 & 26.39 & 2.52$\times$ \\
    QMCPack                    &  68.89 &  61.36 &  2.99 & 29.43 &  2.97 & 2.14$\times$ \\
    Magnetic Reconnection      &  53.88 &  10.96 &  1.96 &  3.89 &  1.95 & 3.44$\times$ \\
    \bottomrule
  \end{tabular}
\end{table}

\paragraph{Results.} Four GPUs reach the quality the single GPU attains
within its training budget $2.1$--$4.1\times$ sooner, cutting compression of
the $512^3$ fields from about $11$--$12$\,s to $3.4$--$5.3$\,s, QMCPack from
$61$--$62$\,s to $23$--$29$\,s, and Miranda from $87$--$95$\,s to
$35$--$39$\,s. Several ratios exceed the hardware step-rate gain alone -- one
optimizer update costs $14.4$\,ms on one GPU and $7.0$\,ms on four, a
$2.1\times$ step-rate ratio, so about half of each step is synchronization --
because reaching equal quality compounds two effects: the four-GPU run steps
faster, and its shard-stratified updates (one slice from each quarter of the
volume) make slightly more progress per update than four independent draws.
That second effect is largest early in training, which is where the
budget-anchored targets sit, and it is why SPERR--temperature reads
$4.08\times$: its crossing lies on the steepest part of both curves. The
serial codec stage adds a further $1$--$15$\,s that no GPU count removes.

Two rows needed the sampler fixed before this comparison was meaningful.
With sequential sampling -- a batch of four \emph{consecutive} slices -- the
single-GPU baseline stalled on Miranda (held at its final quality for
essentially the whole run) and dipped $3$\,dB mid-run on Magnetic
Reconnection, inflating the apparent speed-ups to $46.6\times$ and
$12.4\times$. Both datasets therefore draw slices randomly, as QMCPack always
has, which restored monotone single-GPU curves and the ratios reported above;
Miranda additionally uses the Phase-1 search's choice (axis~1, lr
$3.1\times10^{-3}$), which removed a loss spike that lr $8\times10^{-3}$
produced at this model size. Epoch counts are set so every run covers its
side of the measurement -- QMCPack's single-GPU run gets 20 epochs to fill
its $60$\,s budget, and its four-GPU runs train until they reach the target
-- with the learning-rate schedule annealed over whatever count a run is
given.

Quality is equal within a row by construction, and decompression does not
change: it is a single forward pass on one GPU either way, and the paired
columns agree to within $0.11$\,s (SZ3) and $2.5$\,s (SPERR--Miranda,
dominated by codec decode variance). Decompression costs $1.7$--$29$\,s
against $11$--$95$\,s of single-GPU compression, which suits a
write-once/read-many archive: training is paid once and every later read is
cheap.

Figure~\ref{fig:psnr-vs-time} traces one row per codec in full, showing where
the table's numbers come from. Each curve is cut where it first reaches the
row's target, so the horizontal gap is the measurement; the times shown are
the training share, i.e.\ the table's compression columns minus codec
encoding. On SZ3--NYX the single GPU holds $121.67$\,dB when its $10$\,s
budget expires and the four-GPU run is already there at $2.9$\,s
($3.5\times$); on SPERR--Miranda the $80$\,s budget buys $49.24$\,dB, which
four GPUs reach at $31.7$\,s ($2.5\times$).

\begin{figure}[t]
  \centering
  \includegraphics[width=\columnwidth]{psnr_vs_walltime_2panel_cr500.pdf}
  \caption{Global PSNR against pure training wall time at
  $\mathrm{CR}\!\approx\!500$, from the same runs as
  Table~\ref{tab:parallel-time}: NYX baryon density over SZ3 (left) and
  Miranda over SPERR (right). Within a panel both configurations perform the
  same number of updates per epoch; each curve is cut where it first reaches
  the target quality (dashed line) -- what the single GPU attains within its
  training budget -- so the curves end at the same quality and differ only in
  wall time ($10.0$ vs.\ $2.9$\,s; $80.0$ vs.\ $31.7$\,s). Filled endpoints are
  interpolated crossings, not epoch boundaries; adding codec encoding to them
  gives the compression columns of Table~\ref{tab:parallel-time}.}
  \label{fig:psnr-vs-time}
\end{figure}

% ---------------------------------------------------------------------------
% NOTES FOR REVISION -- not for inclusion
%
% 1. THE SAMPLER HANDICAP -- RESOLVED for Miranda and Magnetic (2026-08-19).
%    With bg_sample_mode "sequential" a batch of n is n CONSECUTIVE slices,
%    whereas four GPUs draw one slice from each of four widely separated
%    shards. Historical evidence of the gap it produced:
%      QMCPack, sequential batch 256 -> 1.7 dB BELOW the SZ3 base
%      Miranda, 238k corrector       -> 4 GPU better by 1.38 dB
%      NYX dark matter, BO config    -> 4 GPU better by 0.94 dB
%    Miranda and Magnetic now use sample="random" (bench_compress_table.py
%    DATASETS), like QMCPack always has, and were rerun into table_rand /
%    table_rand_bo / table_sperr_rand. Effect: monotone 1-GPU curves, 1-GPU
%    finals up (Miranda 43.63 -> 44.99 with BO, Magnetic 50.14 -> 50.66), and
%    the two absurd iso-PSNR ratios collapsed (46.55x -> 4.18x, 12.39x ->
%    2.96x). NYX still trains sequentially -- its rows were never distorted
%    (finals match the 4-GPU runs to within noise) -- but switching it too
%    would make the config uniform; cheap to do if a reviewer asks.
%
% 2. RUN-TO-RUN VARIANCE is up to 0.19 dB for an identical configuration
%    (cudnn.benchmark autotuning + bf16 reduction order). Five of the six
%    quality margins are at or below that scale. Add a footnote stating the
%    figure, or report 3-seed medians.
%
% 3. CUT FROM AN EARLIER DRAFT, still supported by data on disk if a reviewer
%    asks: (a) the per-update cost decomposition, 14.4 ms on one GPU against
%    7.0 ms on four, i.e. 48% of each step is synchronization; (b) the model
%    capacity ablation, Miranda at 54k parameters gives 3.46x against 3.68x at
%    238k, so a 4.4x larger model buys 6% more speed-up while volume takes it
%    from 2.1x at 512^3 to 3.5x at 1024^3 (bench_out/table vs table_240k);
%    (c) the normalization pitfall -- per-shard rather than full-volume stats
%    make the four ranks regress residuals in different unit systems, the
%    per-shard input standard deviations differ by 16x (2.8 to 45.6), and the
%    four-GPU run then ends 0.8 dB BELOW the single-GPU result instead of
%    1.4 dB above it. (c) is the one worth restoring if space allows: it is a
%    silent failure that looks like a large-batch generalization gap.
%
% 4. THE TABLE IS NOW BUDGET-ANCHORED ISO-PSNR (author decision, 2026-08-20),
%    superseding both the min-final iso-PSNR (2026-08-19) and the iso-work
%    framings. Definition: per row, target = the 1-GPU best-so-far PSNR when
%    its training budget expires (TRAIN_BUDGET in bench_compress_table.py:
%    10 s for 512^3 fields, 60 s QMCPack, 80 s Miranda -- matching the main
%    results table's Train column), interpolated on the best-so-far curve,
%    never extrapolated; t1 = the budget, t4 = when the 4-GPU run first
%    reaches the target. Computed by --iso budget (now the default; --iso
%    psnr gives the min-final variant, --iso work the 10-epoch one).
%    CONVENTION CHOICES a reviewer may probe:
%    (i)  The full budget is treated as pure training time. The main table's
%         footnote splits the budget 10%/90% between Phase 1 (BO) and
%         training; under that reading the anchors would be the PSNR at
%         9/54/72 s. Not applied here -- decide once and state it.
%    (ii) Epoch counts differ where the budget requires it: QMC 1-GPU runs 20
%         epochs (10 epochs = 33 s < 60 s budget), QMC 4-GPU runs 20 (SZ3) /
%         40 (SPERR) epochs because the 20-epoch cosine schedule saturated
%         0.13 dB under the target -- schedule-bound, not a DDP quality gap.
%         Everything else runs 10 epochs.
%    KNOWN WEAKNESSES, kept deliberately visible rather than hidden:
%    (a) RESOLVED: the two SZ3 rows that exceeded the 4-GPU bound under
%        sequential sampling (Miranda 46.55x, Magnetic 12.39x) were rerun with
%        random sampling (note 1); they now read 4.18x and 2.96x. The one
%        remaining dagger, Miranda's 4.18x, is a flat-tail artifact: both
%        curves monotone and nearly congruent, equal-work ratio 3.75x, last
%        two 1-GPU epochs gain <0.01 dB each, run-to-run noise 0.19 dB. The
%        caption says so.
%    (b) Time-to-target inherits the 1-GPU curve's shape, which is noisy
%        across settings (earlier evidence: 2.75x vs 5.39x on NYX dark matter,
%        6.35x vs 2.75x on Magnetic between two hyper-parameter sets). The
%        3-seed-median recommendation therefore now applies to the TABLE, not
%        just the figure -- do that before camera-ready if possible. Miranda's
%        marginal 4.18x is exactly the kind of number a 3-seed median would
%        settle.
%    (c) The 4-GPU quality bonus at full budget (e.g. 123.12 vs 122.94 on NYX
%        baryon) no longer appears in the table; it survives as one prose
%        sentence. Reviewers comparing against the figure's endpoints will see
%        consistent numbers by construction.
%    Figure endpoints and table columns now agree: figure times are the
%    table's Comp. minus the codec share.
%
% 5. CURVE WOBBLE IS NOISE, NOT LEARNING RATE. Single-GPU curves sometimes dip
%    several dB mid-run and recover. Re-running under Phase-1 BO, which chose an
%    even higher lr, moved the wobble from 1.91 dB to 0.00 dB on NYX baryon
%    density but from 0.07 dB to 5.34 dB on NYX dark matter. Four paired runs
%    show no consistent direction.
%
% 6. HYPER-PARAMETERS ARE FIXED (axis 0, lr 8e-3), not searched per run, so the
%    GPU count is not confounded with tuning. A control repeating all six
%    datasets under the Phase-1 (axis, lr) search changed the speed-up by at
%    most 4.5% (bench_out/table_bo). Two limits are worth stating if the paper
%    claims Phase-1 helps: (a) the downsampled proxy cannot separate close
%    candidates -- on NYX dark matter it preferred axis 2 over axis 0 by 0.01 dB
%    on the proxy and the full run came out 0.90 dB WORSE; (b) with 10 trials
%    and one enqueued per axis, TPE spends the rest on the leading axis, so the
%    per-axis numbers in its log are not a fair axis comparison.
%
% 7. bg_h WAS CAPPED, NOT BUDGETED. pick_bg_h_under_budget defaults to
%    h_candidates = range(3, 30), so Miranda's 240k budget selected h=29 =
%    54,466 parameters, i.e. 22% of it. bench_compress_table.py now passes
%    range(3, 128), and Table~\ref{tab:parallel-config} reports the resulting
%    237,538-parameter model.
%
% 8. NOT YET MEASURED: scaling past 4 GPUs; multi-node.
%
% 9. SPERR-MIRANDA RUNS AT lr 4e-3, NOT THE GLOBAL 8e-3 (the random-sampling
%    rerun in table_sperr_rand keeps 4e-3; SZ3-Miranda now runs at the
%    Phase-1 lr 3.1e-3, so neither Miranda row is at 8e-3 -- consistent with
%    this model size being unstable there). At 8e-3 the 4-GPU DDP
%    run is unstable on SPERR's Miranda residual: seed 17 never rose above the
%    46.80 dB base (ended 46.79) and seed 18 diverged to 43.74, while the
%    1-GPU run reached 48.23. Halving the lr fixes it for BOTH columns (49.33
%    1-GPU / 49.31 4-GPU), so the row is internally iso-config and fair; it
%    just deviates from the "one lr everywhere" statement in note 6. The
%    unstable 8e-3 pair is kept in bench_out/probe/miranda_n*_lr8e3_unstable
%    .json. If a reviewer asks why SPERR's envelope matters: SPERR targets
%    PSNR, not L-inf, so its Miranda rel envelope is 8.7e-2 against SZ3's
%    1.6e-2 -- a ~5x looser clamp, a larger residual target, and a harder
%    optimization at high lr.
%
% 10. SPERR IS TIMED THROUGH ITS CLI (sperr3d subprocess: process start + file
%    read + bitstream write), SZ3 in-process via pysz on an in-memory array.
%    Each codec block is internally consistent, but the ~1s SPERR codec share
%    on 512^3 fields is not comparable to SZ3's 1.6s share. The speed-up
%    column is training-only, so it is unaffected.
% ---------------------------------------------------------------------------
